% Figure for Thermal Physics §A4.3 (proof of the equivalence of the Clausius and Kelvin-Planck statements: combining an impossible device with an ordinary one; after University Physics Vol. 2, Fig. 4.9, and PHY 287 Oct 13 slides).
\begin{tikzpicture}[font=\small, >={Stealth[length=2.2mm]},
    res/.style={draw=black!70, thick, fill=black!6, minimum height=0.55cm, font=\scriptsize},
    dev/.style={draw=#1, very thick, fill=#1!12, circle, minimum size=1.15cm, inner sep=1pt, font=\scriptsize, align=center},
    fl/.style={->, very thick, #1}, lab/.style={font=\scriptsize, fill=white, inner sep=1pt}]
  % ================= row 1: Clausius false => Kelvin-Planck false =================
  \begin{scope}
    \node[font=\scriptsize, black!70, anchor=west] at (-0.9,3.05) {Clausius false $\Rightarrow$ Kelvin--Planck false};
    \node[res, minimum width=5.0cm] (H) at (1.6,2.3) {hot reservoir $T_h$};
    \node[res, minimum width=5.0cm] (C) at (1.6,-1.3) {cold reservoir $T_c$};
    \node[dev=figR] (P) at (0.2,0.5) {perfect\\refrig.};
    \node[dev=figB] (E) at (2.9,0.5) {engine};
    \draw[fl=figR] (0.2,-1.0) -- (P.south) node[midway, lab, left=2pt] {$Q$};
    \draw[fl=figR] (P.north) -- (0.2,2.0) node[midway, lab, left=2pt] {$Q$};
    \draw[fl=figB] (2.9,2.0) -- (E.north) node[midway, lab, right=2pt] {$Q + \Delta Q$};
    \draw[fl=figB] (E.south) -- (2.9,-1.0) node[midway, lab, right=2pt] {$Q$};
    \draw[fl=figB] (E.east) -- ++(0.8,0) node[right, font=\scriptsize] {$W = \Delta Q$};
    \node[font=\large] at (6.4,0.5) {$=$};
    \node[res, minimum width=2.2cm] (H2) at (8.4,2.3) {hot $T_h$};
    \node[res, minimum width=2.2cm] (C2) at (8.4,-1.3) {cold $T_c$};
    \node[dev=figR] (N) at (8.4,0.5) {perfect\\engine};
    \draw[fl=figR] (8.4,2.0) -- (N.north) node[midway, lab, right=2pt] {$\Delta Q$};
    \draw[fl=figR] (N.east) -- ++(0.8,0) node[right, font=\scriptsize] {$W$};
    \node[font=\scriptsize, black!60] at (8.4,-0.55) {no net heat};
  \end{scope}
  % ================= row 2: Kelvin-Planck false => Clausius false =================
  \begin{scope}[yshift=-5.0cm]
    \node[font=\scriptsize, black!70, anchor=west] at (-0.9,3.05) {Kelvin--Planck false $\Rightarrow$ Clausius false};
    \node[res, minimum width=5.0cm] (H) at (1.6,2.3) {hot reservoir $T_h$};
    \node[res, minimum width=5.0cm] (C) at (1.6,-1.3) {cold reservoir $T_c$};
    \node[dev=figR] (P) at (0.2,0.5) {perfect\\engine};
    \node[dev=figB] (R) at (2.9,0.5) {refrig.};
    \draw[fl=figR] (0.2,2.0) -- (P.north) node[midway, lab, left=2pt] {$Q$};
    \draw[fl=figR] (P.east) -- (R.west) node[midway, lab, above=2pt] {$W = Q$};
    \draw[fl=figB] (2.9,-1.0) -- (R.south) node[midway, lab, right=2pt] {$Q_c$};
    \draw[fl=figB] (R.north) -- (2.9,2.0) node[midway, lab, right=2pt] {$Q_c + W$};
    \node[font=\large] at (6.4,0.5) {$=$};
    \node[res, minimum width=2.2cm] (H2) at (8.4,2.3) {hot $T_h$};
    \node[res, minimum width=2.2cm] (C2) at (8.4,-1.3) {cold $T_c$};
    \node[dev=figR] (N) at (8.4,0.5) {perfect\\refrig.};
    \draw[fl=figR] (8.4,-1.0) -- (N.south) node[midway, lab, right=2pt] {$Q_c$};
    \draw[fl=figR] (N.north) -- (8.4,2.0) node[midway, lab, right=2pt] {$Q_c$};
    \node[font=\scriptsize, black!60, anchor=west] at (9.1,0.5) {no work};
  \end{scope}
\end{tikzpicture}
